% Introduction: background, research gap, hypotheses, contributions, structure.
% Author-year references are not used here; the theoretical citations underpinning
% this chapter are developed in full in Section~\ref{sec:literature-review}.
% Restructured 21 August 2026 to foreground H2 and the extensive-/intensive-margin
% decomposition as the primary contribution and to move MI/HA/WH detail to
% Sections~\ref{sec:research-design}-\ref{sec:results}. No hypothesis decision,
% coefficient, or sample size was changed.

\section{Introduction}
\label{sec:introduction}

Prediction markets aggregate dispersed beliefs into continuously updated contract prices that are commonly interpreted as market-implied probabilities, subject to assumptions about preferences and market structure. Evidence from political and other event markets shows that these prices can provide useful forecasts, although their calibration and informational content vary across settings. Large individual trades are of particular interest within this setting because market microstructure theory predicts that trade size carries information about who is trading and why, yet whether an unusually large trade on a continuously repriced, retail-accessible platform behaves like an informed order, a liquidity-consuming order, or a source of temporary price pressure is an open empirical question that this dissertation is designed to inform rather than resolve.

Polymarket provides an unusual setting in which to study this question. It combines off-chain order matching with blockchain settlement and publicly observable trading records, so trade size, direction, and timing can be measured directly and reproducibly from public data. Public records do not, however, reveal beneficial ownership, private information, or trading intent, so this dissertation's empirical question is narrower than a direct test of informed trading or misconduct: how are unusually large trades associated with subsequent probability adjustment, and does that association more plausibly reflect information, mechanical execution, or strategic trading? Throughout this dissertation, ``price impact'' denotes a reduced-form short-horizon conditional association between signed trading flow and subsequent prices, not a causal structural price-impact estimate; Section~\ref{sec:research-design} states this definition precisely, and every subsequent chapter applies it consistently.

A further conceptual distinction motivates the design. Existing discussion of large Polymarket participants, in both the trading community and the emerging 2025--2026 academic literature, has largely emphasised wallet-level activity, profitability, or skill. This framing conflates three questions that this dissertation keeps separate: whether an individual trade is unusually large relative to its own market, whether an observed wallet is unusually active over time, and whether an observed wallet has deployed unusually large cumulative capital. Recent cross-platform evidence classifies individual trades by relative size and studies their aggregated directional imbalance, but evidence across a broad set of sports contracts remains limited, and no existing study jointly separates trade size from wallet activity, decomposes the resulting price association into whether an update occurs versus how large it is conditional on occurring, and tests out-of-sample predictability within one design. This dissertation addresses that setting-specific gap using 6,725,732 canonical trades across 360 Polymarket markets linked to the 2026 FIFA World Cup: 312 match contracts nested in 104 matches and 48 country-outright contracts.

The primary empirical question is whether flow generated by unusually large individual trades, defined using a market-specific past-only 99th-percentile value threshold, has a stronger short-horizon price association than flow from ordinary-sized trades (H2), after first confirming that signed total flow is associated with short-horizon price change at all (H1). An extensive-/intensive-margin decomposition is treated as a mechanism diagnostic rather than embedded in the primary hypothesis. Two further hypotheses and a research question ask whether large-trade-associated movements persist and improve forecast accuracy, or instead reverse (H3a and H3b), and how the association varies with time to kickoff and match phase (RQ4). Supplementary market-integrity analysis and secondary wallet-level classifications separate trade size from wallet activity and capital deployment.

Three findings anchor the completed analysis. First, under both expanding-history and rolling seven-day past-only P99 definitions, large-trade flow has a larger five-minute directional association than ordinary flow. The estimated large-minus-ordinary contrasts are 0.000082 and 0.000081, respectively. Second, placebo lead estimates are indistinguishable from zero, whereas the association appears after the trade at one, five, and fifteen minutes and survives event-by-time fixed effects. This timing pattern is consistent with short-horizon price response, although it does not establish causality or identify the underlying mechanism. Third, the extensive-margin interpretation is model-sensitive: linear probability models find a larger association of large flow with the occurrence of a price update, but market-fixed-effects logit models do not find a statistically distinguishable large-versus-ordinary difference. Persistence, forecast-accuracy, behavioural, predictive, and anomaly-screening analyses therefore provide qualifying evidence rather than proof of informed trading or manipulation.

This dissertation contributes academically by implementing past-only classifications that avoid look-ahead, separating trade-level size from wallet-level constructs, testing temporal ordering with lead--lag estimates, evaluating identification under event-by-time fixed effects, and showing that conclusions can depend on functional form. It also separates explanatory fixed-effects evidence from genuinely out-of-sample prediction and unsupervised anomaly discovery. The practical contribution is a reproducible, non-causal framework for studying unusual trading behaviour without treating trade size, wallet concentration, predictive performance, or an anomaly score as evidence of misconduct.

The remainder of the dissertation reviews the relevant theory and evidence,
documents the data and research design, reports the primary and supplementary
results, and then discusses their implications and limitations before
concluding.
